\subsection*{Published window normalization and an external numerical lock}

The support half-width used by Zhu~\cite{Zhu2026Windows} is
$a=\log\lambda$; the present Fourier interval has full length
$L=2a$. These symbols must not be interchanged. The comparison below
uses the same physical Weil form and ordinary $L^2$ norm, after
translation of the logarithmic interval.
\begin{center}
\begin{tabular}{lrr}
\toprule
Result & half-width $a$ & multiplicative $\lambda=e^a$\\
\midrule
Classical Yoshida / Connes--Consani & $\tfrac12\log2\simeq0.3466$ & $\sqrt2$\\
Zhu's certified window & $0.8$ & $e^{0.8}\simeq2.2255$\\
Proposition~\ref{prop:v116-window-positive} & $\log3\simeq1.0986$ & $3$\\
Proposition~\ref{prop:v126-full-window} & $\log4\simeq1.3863$ & $4$\\
\bottomrule
\end{tabular}
\end{center}
Thus the certified $\lambda=4$ half-width is approximately $1.73$
times Zhu's and exactly four times the classical half-width. Zhu
provides a quantitative lower ground bound on the smaller window;
Proposition~\ref{prop:v140-ground4} supplies the complete simple-even
ground ordering here. No ordinary ground lower constant may be read
from an LDL pivot or a relative Schur margin. The empirical
Landau--Widom fit is not a normalization test: unconverged finite
compression values are only upper bounds for the complete infimum.

In the notation of Connes--Consani--Moscovici
\cite[Corollaries 3.7--3.8]{ConnesConsaniMoscovici2025},
$\mu^{\mathrm{CCM}}_\lambda=\inf\sigma(A_\lambda)$ decreases with $\lambda$.
Support consistency and the common form normalization therefore give
\[
 \mathsf W_4\succeq0\quad\Longleftrightarrow\quad
 \mu^{\mathrm{CCM}}_\lambda\ge0\ \text{for every }1<\lambda\le4.
\]
The fixed floors in Proposition~\ref{prop:v138-three-floors} similarly
bound the corresponding complete $\mu^{\mathrm{CCM}}_\lambda$. These are fixed-range
statements. A cofinal bound is an additional open input.

\begin{lemma}[The centering phase in the Fourier dictionary]
\label{lem:v141-centering}
Put $a=L/2$, $\omega_n=2\pi n/L$, and let $v=(v_0,\ldots,v_N)$
be coefficients in the actual unshifted even basis on $[0,L]$:
\[
 f_v(y)=\frac{v_0}{\sqrt L}
       +\sqrt{\frac2L}\sum_{n=1}^Nv_n\cos(\omega_n y).
\]
Its centered representative $g_v(x)=f_v(x+a)$ on $[-a,a]$ has
cosine coefficients $(-1)^nv_n$ for $n\ge1$. For the nonunitary
Fourier transform used by CCM,
\begin{equation}\label{eq:v141-transform}
 \widehat g_v(z)=\frac{2\sin(az)}{\sqrt L}
 \left\{\frac{v_0}{z}+\sqrt2\sum_{n=1}^N
          \frac{v_nz}{z^2-\omega_n^2}\right\}.
\end{equation}
All apparent poles are removable. In particular, for $1\le k\le N$,
\begin{equation}\label{eq:v141-lattice-value}
 \widehat g_v(\omega_k)=\sqrt{L/2}\,(-1)^kv_k.
\end{equation}
Hence the sine factor alone does not impose lattice zeros in the
retained band.
\end{lemma}
\begin{proof}
The phase is the exact identity
$\cos(\omega_n(x+a))=(-1)^n\cos(\omega_nx)$.
Integrating each cosine against $e^{-izx}$ cancels this phase
against the factor $(-1)^n$ in the sine numerators, giving
\eqref{eq:v141-transform}. Orthogonality of exponentials, or the
removable-pole limit, gives \eqref{eq:v141-lattice-value}.
This is CCM's finite partial-fraction formula in orthonormal parity
coordinates. Unitary Fourier normalization multiplies it by
$(2\pi)^{-1/2}$; translating the whole function multiplies its
transform by a zero-free exponential. Neither changes its zeros.
\end{proof}

\begin{proposition}[A certified finite-compression semantic lock]
\label{prop:v141-ccm-lock}
At the already studied window $\lambda=3$, take $N=120$ as in
\cite[Section 6, Figure 1]{ConnesConsaniMoscovici2025}. Assemble
both parity compressions with the existing exact coefficient formulas.
Their common global lowest eigenvalue is simple and even. Normalize
its even coefficient vector by $v_0=1$. For each of the first eight
positive zeta ordinates $\gamma_k$, the transform
\eqref{eq:v141-transform} has a unique root in the archived local
interval $I_k$ of radius at most $1.000001\cdot10^{-110}$.
Its positive discrepancy from $\gamma_k$ rounds to the corresponding
published entry:
\begin{center}
\begin{tabular}{rrr}
\toprule
$k$ & reproduced discrepancy (rounded) & CCM Figure 1\\
\midrule
1 & $1.582329697193127\cdot10^{-34}$ & $1.6\cdot10^{-34}$\\
2 & $2.060550289594801\cdot10^{-31}$ & $2.1\cdot10^{-31}$\\
3 & $1.460904856368295\cdot10^{-29}$ & $1.5\cdot10^{-29}$\\
4 & $8.293686860378121\cdot10^{-27}$ & $8.3\cdot10^{-27}$\\
5 & $1.279455965166440\cdot10^{-25}$ & $1.3\cdot10^{-25}$\\
6 & $1.174389453606719\cdot10^{-23}$ & $1.2\cdot10^{-23}$\\
7 & $7.528486586592898\cdot10^{-22}$ & $7.5\cdot10^{-22}$\\
8 & $6.565874264666424\cdot10^{-21}$ & $6.6\cdot10^{-21}$\\
\bottomrule
\end{tabular}
\end{center}
This is a local finite-compression root certificate, not an exhaustive
ordering of all transform zeros or a transfer to the complete ground.

\paragraph{Normalization scope of the numerical test.}
For any $c>0$ and real $s$, replacing the assembled finite matrix
$A$ by $cA+sI$ preserves its ground eigenspace and hence these
transform zeros. Agreement with the published zero discrepancies
cannot by itself test an overall positive scalar or an identity shift.
The identification of the full form normalization rests separately on
the defining coefficient formulas, not on this numerical benchmark.
\end{proposition}
\begin{proof}[Interval verification]
The verifier in \texttt{evidence/v141/} uses the original archived
coefficient assembler, \texttt{assembly\_general.py}, without modification.
It reassembles $121$ even and $120$ odd coordinates using the
$128$-term archimedean formula and its explicit geometric remainder.
The logarithmic diagonal, normalized zero mode, pole term and all
prime powers below $9$ are retained. The endpoint contribution at
$9$ vanishes on the autocorrelations, so the strict prime cutoff
agrees with CCM's inclusive convention.

Arb's verified eigenpair routine isolates all eigenvalues and proves
that the chosen even value lies strictly below every other even and
odd value. Verified eigenvector ratios enclose the real coefficients.
Midpoint root finding only proposes decimal centers. Acceptance uses
opposite interval signs of the rational factor at the endpoints,
a derivative interval excluding zero on the entire bracket, and a
sine interval excluding zero there. These give existence and local
uniqueness by the intermediate value theorem and monotonicity.
Arb enclosures of the fixed low zeta ordinates then give the saved
discrepancy intervals, each strictly inside its published rounding
cell. Every gate is replayed at $768$ and $1024$ bits, at the same
$\lambda,N$ and series order. Direct termwise transforms provide an
additional algebraic consistency check. No midpoint solve is used.
\end{proof}

\paragraph{What had to change.}
The earlier NS-14 diagnostic applied unshifted eigenvector coefficients
to centered cosines without the required $(-1)^n$ conversion. It
therefore evaluated a different window function. In addition, a small
coefficient is not zero: \eqref{eq:v141-lattice-value} precludes the
inference from concentrated low-mode mass to exact retained-band
lattice zeros. The new certificate disproves that diagnostic inference
at $\lambda=3$ and reproduces CCM's fixed-$N$ benchmark. It does not
establish complete-ground zero locations at $\lambda=3$ or $4$, nor
convergence to $\Xi$. Those transfer and cofinal inputs remain open;
G2 and RH remain open. No new window is computed.
